% !PRINTED 739
\weberpage{739}

%<<CONT
then from equations (13) we obtain
\[
2 \sum\limits_{1,\, \nu-1}^{i} a_i\, l_{i,k} = \xi_0 L_{0,k} + \xi_1 L_{1,k} + \cdots + \xi_{\nu'-1} L_{\nu'-1,k}.
\]
From this there results, by (10),
\[
\sum^{i} a_i\, l_{i,k+\nu'} = -\sum^{i} a_i\, l_{i,k},
\]
and this means that
\[
\delta = \varepsilon_1^{a_1}\, \varepsilon_2^{a_2} \ldots \varepsilon_{\nu-1}^{a_{\nu-1}}
\]
is a normal unit of the field $H_m$. Since $\mathcal{E}_0, \mathcal{E}_1, \ldots, \mathcal{E}_{\nu'-1}$
is a fundamental system of normal units, it then follows further that
 the numbers $\xi_0, \xi_1, \ldots, \xi_{\nu'-1}$ must be even. From this
there results:

1. \gesperrt{The system of equations (18) has the congruences
}
\weq{19}{\begin{gathered}
\sum\limits_{1,\, \nu-1}^{i} a_i\, M_{k,i} \equiv 0 \pmod{2} \\
k = 0, 1, \ldots, \nu'-1
\end{gathered}}
\gesperrt{as a consequence.}

But equations (18), whose number amounts to only $\nu'-1$, contain
$\nu-1$ unknowns $a_i$. Hence there are several systems of mutually
independent solutions, which we can make somewhat more precise
as follows:

2. \gesperrt{There are $\nu'$ integral solutions of the system
(18):}
\weq{20}{
\begin{aligned}
& a_{1,1}, a_{2,1}, \ldots, a_{\nu-1,1} \\
& a_{1,2}, a_{2,2}, \ldots, a_{\nu-1,2} \\
& \ldots\ldots\ldots\ldots \\
& a_{1,\nu'}, a_{2,\nu'}, \ldots, a_{\nu-1,\nu'}
\end{aligned}
}
\gesperrt{of such a kind that from the matrix (20) a
determinant of $\nu'$ rows can be formed
which is an odd number.}

For there is in the first place a solution of (18):
\[
a_{1,1}, a_{2,1}, a_{3,1}, \ldots, a_{\nu-1,1},
\]
in which an odd number occurs. Let us assume---as is permissible,
since here nothing depends on the arrangement of the indices---that
%CONT>>
